\chapter{Factor V Leiden Test}

\section*{What Is This Test?}
Factor V Leiden testing detects the inherited factor V variant that makes factor V relatively resistant to activated protein C, an important natural anticoagulant. The usual cause is the c.1601G>A (p.Arg534Gln) variant, historically called R506Q. It increases the tendency toward venous thromboembolism, particularly when combined with acquired risks such as surgery, estrogen exposure, pregnancy, immobility, or cancer.

\section*{Alternative Names}
Factor V Leiden Mutation Test; F5 c.1601G>A; Factor V R506Q; Activated Protein C Resistance Genetic Test; Thrombophilia Genetic Test

\section*{How It Is Performed}
The genetic test is usually performed on EDTA-anticoagulated blood using PCR-based allelic discrimination or another validated molecular method. An activated-protein-C-resistance clotting assay may be used as an initial functional screen, but genetic testing is useful for confirmation and is less affected by anticoagulant treatment or pregnancy.

\section*{Why It Is Ordered}
\begin{itemize}
  \item Selected evaluation of unexplained or recurrent venous thromboembolism
  \item Assessment of a strong or unusual family history of venous thrombosis
  \item Clarification of an abnormal activated-protein-C-resistance result
  \item Thrombophilia assessment when the result would change counseling or management
\end{itemize}

\section*{Interpretation and Limitations}
A heterozygous result indicates one affected copy of F5; homozygosity indicates two. A positive result identifies a risk factor but does not prove that a clot was caused by the variant or predict exactly when a clot will occur. Testing is not routinely needed for every patient with thrombosis, and a negative result does not exclude other inherited or acquired thrombophilias. Interpretation should be guided by the clinical context and current thrombosis guidance.
\section*{Specimen and Preparation}
The usual specimen is EDTA-anticoagulated whole blood; fasting is not required. Tell the clinician about pregnancy, estrogen-containing medicines, anticoagulants, recent transfusion, and previous thrombosis. These factors may affect a clot-based activated-protein-C-resistance screen, whereas the molecular test measures the inherited DNA variant. Genetic counseling is useful before testing because results can have implications for relatives.

\section*{Risks and Safety}
Physical risk is limited to brief pain, bruising, or bleeding from venipuncture. A positive result identifies a risk factor but does not prove that a clot was caused by it or require lifelong anticoagulation. New chest pain, breathlessness, coughing blood, or one-sided leg swelling requires urgent assessment regardless of the test result.

\section*{Understanding Results and Follow-up}
A negative result means the tested Factor V Leiden variant was not detected; it does not exclude other inherited or acquired thrombophilias. Follow-up may include review of personal and family history, counseling of selected adult relatives, and evaluation for other causes when clinically justified. Do not start, stop, or change anticoagulant or hormone treatment solely from this result.

\section*{Regulatory Status and Authorization Date}
This chapter describes a test category, not a specific commercial product. FDA, CDSCO, EU IVDR, and MHRA approval or registration dates therefore cannot be assigned without the assay manufacturer, product name, instrument, and intended use. For laboratory-developed tests, an individual agency approval date may not exist. See the Preface for the interpretation of regulatory status.

\clearpage
